\section{Evaluation}\label{s:evaluation}

This section evaluates the correctness, security, and performance of the proposed system. The experiments are designed to verify that the protocol behaves as expected under both normal and adverse conditions, while also assessing its scalability and sensitivity to physical-layer parameters.

For each experiment, we define the simulation setup, the metrics collected, and the expected outcome. The primary metrics considered throughout this evaluation are the CHSH \(S\) parameter, key correctness, simulation execution time, and protocol termination behavior. 

To evaluate the performance and correctness of the system, we will conduct the following tests:

\begin{itemize}
  \item \textbf{1. Simulation run with two client nodes and an EPR source}: This is the simplest and most important test. For a successful run, the CHSH \(S\) parameter should be greater than 2, indicating a violation of the classical bound, and the keys generated by both nodes must be identical.

  \item \textbf{2. Simulation run over an insecure channel}: This test demonstrates the effects of operating the simulation over an insecure channel. For a successful test, the CHSH \(S\) parameter should fall below 2, indicating the absence of quantum correlations, and the protocol should terminate the key-generation process.

  \item \textbf{3. Relationship between CHSH values and distance}: This test examines the impact of distance on fidelity degradation and, consequently, on the CHSH value. Since fidelity degradation is expected to increase with distance, the measured CHSH value should decrease accordingly.

  \item \textbf{4. Scalability of the system}: This test evaluates the system's ability to handle simulations of varying sizes, ranging from small-scale runs involving 100 qubits to large-scale runs involving 100{,}000 qubits.

  \item \textbf{5. Impact of dark counts on simulation time}: This test investigates how detector dark counts affect detector cooldown periods and, consequently, the total time required to complete a simulation run. An increase in the dark count rate is expected to increase the overall simulation time.
\end{itemize}

\subsection{Simulation run with two client nodes and an EPR source}

This experiment validates the basic functionality of the system under normal operating conditions. The objective is to verify that entanglement distribution, CHSH-based security verification, and key generation are performed correctly.

% TODO: MAKE PICTURE OF TOPOLOGY ONCE WEB IS FINISHED
The simulation consists of two client nodes, denoted as \texttt{src} and \texttt{dst}, connected to a shared EPR source through secure links. Each client is connected to the EPR node via a channel of 100\,000 meters (100 km). Figure~\ref{...} illustrates the network topology used in this experiment.

For the protocol to be considered successful, the measured CHSH \(S\) parameter must be greater than 2, demonstrating the presence of quantum correlations, and both client nodes must generate identical secret keys.

In order to generate a final secret key of 1024 bits, the protocol must distribute at least 10,000 entangled pairs. Since the CHSH protocol employs three possible measurement bases per participant, only approximately $\frac{1}{9}$ of the distributed pairs are expected to be measured using the same basis combination required for key generation.

As the CHSH S-paremeter is not deterministic, we will perform 10 runs and average the CHSH value accross all of them. 

The results of the 10 runs are shown in Table~\ref{tab:chsh_runs}

\begin{table}[h]
\centering
\caption{CHSH values obtained over 10 simulation runs.}
\label{tab:chsh_runs}
\begin{tabular}{cc}
\toprule
\textbf{Run} & \textbf{CHSH Value} \\
\midrule
1  &  2.488\\
2  &  2.468\\
3  &  2.503\\
4  &  2.452\\
5  &  2.475\\
6  &  2.439\\
7  &  2.459\\
8  &  2.357\\
9  &  2.434\\
10 &  2.410\\
\bottomrule
AVG & 2.449\\
\end{tabular}
\end{table}

As the runs yield an average CHSH value of \(2.449\), this confirms that the correlations between the measurement bases in the test qubits are correct and demonstrates the presence of entanglement between the qubits. We now examine the generated keys.

For brevity, only the keys from a single run are shown, and they are presented in hexadecimal format.

The two keys generated were:
\begin{align*}
\text{Source key:}      &\quad \texttt{f4aa5b5\ldots cb6ba6} \\
\text{Destination key:} &\quad \texttt{f4aa5b5\ldots cb6ba6}
\end{align*}

The results confirm that the system operates correctly. All 10 runs produced a CHSH \(S\) value exceeding the classical bound of 2, with an average of \(2.449\). The deviation from the theoretical maximum is expected, as fidelity degradation over the 100\,km channel introduces noise into the entangled pairs.In addition to this, the source and destination keys are identical across all runs, confirming that the sifting and error correction function as intended

\subsection{Simulation run over an insecure channel}

This experiment validates the basic functionality of insecure channels. The objective is to verify that CHSH-based security verification and mallory actuation are performed correctly.

% TODO: MAKE PICTURE OF TOPOLOGY ONCE WEB IS FINISHED
The simulation consists of two client nodes, denoted as \texttt{src} and \texttt{dst}, connected to a shared EPR source through a secure and insecure link respectively. Each client is connected to the EPR node via a channel of 100\,000 meters (100 km). Figure~\ref{...} illustrates the network topology used in this experiment.

For the protocol to be considered successful, the measured CHSH \(S\) parameter must be below 2, demonstrating the presence of a third person that breaks quantum correlations, and the simulation run must be stopped.

As the CHSH S-paremeter is not deterministic, we will perform 10 runs and average the CHSH value accross all of them. 

The results of the 10 runs are shown in Table~\ref{tab:chsh_runs_insecure}

\begin{table}[h]
\centering
\caption{CHSH values obtained over 10 simulation runs in an insecure channel.}
\label{tab:chsh_runs_insecure}
\begin{tabular}{cc}
\toprule
\textbf{Run} & \textbf{CHSH Value} \\
\midrule
1  &  0.444\\
2  &  0.578\\
3  &  0.517\\
4  &  0.476\\
5  &  0.461\\
6  &  0.559\\
7  &  0.533\\
8  &  0.567\\
9  &  0.399\\
10 &  0.380\\
\bottomrule
AVG & 0.491\\
\end{tabular}
\end{table}

As shown in Table~\ref{tab:chsh_runs_insecure}, the average CHSH value across all runs is \(0.491\). This significant reduction below the classical threshold of \(S = 2\) indicates the absence of quantum correlations in the presence of an adversarial or insecure channel.

Additionally, the greater variance in the measurement outcomes suggests increased randomness, consistent with the disruption of entanglement correlations under insecure conditions.

In addition to this, we can see that the simulator displays the following error:

\begin{quote}
Mallory was detected with CHSH 0.380 in process 0
\end{quote}

The results confirm that the system operates as intended in the presence of insecure channels. All 10 runs produced a CHSH $S$ value below the classical bound of 2. As expected, all runs were halted due to the detection of Mallory.
